\documentclass[11pt,a4paper]{article}

\usepackage[T1]{fontenc}
\usepackage[utf8]{inputenc}
\usepackage{lmodern}
\usepackage[french]{babel}
\usepackage[margin=2.2cm]{geometry}
\usepackage{xcolor}
\usepackage{mdframed}
\usepackage{enumitem}
\usepackage{booktabs}
\usepackage[hidelinks]{hyperref}

\definecolor{primary}{HTML}{1F3B6E}
\definecolor{accent}{HTML}{E26D3C}
\definecolor{yann}{HTML}{1F3B6E}
\definecolor{venus}{HTML}{2E7D5B}
\definecolor{faical}{HTML}{B5482A}
\definecolor{panel}{HTML}{F3F6FB}

\setlength{\parindent}{0pt}
\setlength{\parskip}{0.4em}

% Un tour de parole : étiquette colorée du présentateur, puis le texte à dire.
\newenvironment{turn}[2]%
{\begin{mdframed}[linewidth=0pt,backgroundcolor=panel,skipabove=6pt,skipbelow=6pt,
   innerleftmargin=10pt,innerrightmargin=10pt,innertopmargin=7pt,innerbottommargin=7pt,
   leftline=true,linecolor=#1,linewidth=3pt]%
 {\color{#1}\bfseries #2}\par\vspace{2pt}}
{\end{mdframed}}

\begin{document}

% ----------------------------- Titre ----------------------------------
\begin{center}
{\color{primary}\LARGE\bfseries Script de présentation orale}\\[4pt]
{\large Predicting User Sentiment in Software Issue Reports}\\[10pt]
{\large\bfseries Yann MASSOUAM \quad$\cdot$\quad Vénus BAKIKO \quad$\cdot$\quad Faïçal DIELO}\\[4pt]
{ST2MLE : Machine Learning for IT Engineers}
\end{center}

\vspace{0.4cm}
{\color{accent}\rule{\linewidth}{1.2pt}}

\textbf{Mode d'emploi.} Ce document répartit le pitch entre les trois présentateurs,
slide par slide. Chaque bloc coloré indique qui parle et le texte à dire. Durée visée :
environ 18 à 20 minutes, soit à peu près une minute par slide. Les mêmes textes sont
disponibles dans les notes du présentateur du fichier PowerPoint.

\begin{table}[h]
\centering
\begin{tabular}{lll}
\toprule
\textbf{Présentateur} & \textbf{Slides} & \textbf{Partie} \\
\midrule
\textcolor{yann}{\textbf{Yann MASSOUAM}}  & 1 à 6   & Introduction, données et début de la chaîne \\
\textcolor{venus}{\textbf{Vénus BAKIKO}}  & 7 à 12  & Vectorisation, modèles, métriques et résultats \\
\textcolor{faical}{\textbf{Faïçal DIELO}} & 13 à 18 & Réglages, conclusion et perspectives \\
\bottomrule
\end{tabular}
\end{table}

% ============================ YANN ====================================
\section*{\textcolor{yann}{Partie 1 : Yann MASSOUAM (slides 1 à 6)}}

\begin{turn}{yann}{Slide 1 : Page de titre}
Bonjour à tous. Nous sommes Yann, Vénus et Faïçal, et nous vous présentons notre projet
de Machine Learning : prédire le sentiment d'avis d'utilisateurs sur des applications,
en trois classes, négatif, neutre et positif. En une vingtaine de minutes, nous allons
montrer comment, à partir d'un simple texte d'avis, on construit une chaîne complète qui
nettoie, transforme, entraîne et évalue plusieurs modèles, puis nous expliquerons
pourquoi certains choix comptent vraiment.
\end{turn}

\begin{turn}{yann}{Slide 2 : Plan}
Voici notre fil conducteur. On part du problème, puis des données, puis on déroule la
chaîne de traitement de gauche à droite. On insiste sur deux idées fortes : la
vectorisation, c'est à dire le passage du texte aux nombres, et les métriques, parce que
sur ce projet une bonne exactitude peut cacher un modèle qui ne fonctionne pas. On
termine par les réglages classiques et la conclusion.
\end{turn}

\begin{turn}{yann}{Slide 3 : Le problème}
Le but est concret : à partir du texte d'un avis, prédire automatiquement s'il est
négatif, neutre ou positif. C'est de l'apprentissage supervisé, car chaque avis
d'entraînement possède déjà une étiquette. L'intérêt pratique est réel : sur une
application qui reçoit des milliers d'avis, on veut repérer vite les mécontents, suivre
la satisfaction, et éviter de tout lire à la main.
\end{turn}

\begin{turn}{yann}{Slide 4 : Les données}
Point important : nous n'utilisons pas de données inventées, mais de vrais avis Facebook
téléchargés sur Kaggle. Chaque avis a une note de une à cinq étoiles que l'on transforme
en sentiment : une ou deux étoiles, c'est négatif, trois, c'est neutre, quatre ou cinq,
c'est positif. On supprime les doublons de texte pour éviter qu'un même avis se retrouve
à la fois en entraînement et en test. Il reste 5\,923 avis. Regardez le déséquilibre :
deux tiers d'avis positifs et seulement 4,5 pour cent de neutres. Ce déséquilibre sera au
cœur de notre analyse.
\end{turn}

\begin{turn}{yann}{Slide 5 : La chaîne de traitement}
Voici la chaîne complète, à lire de gauche à droite. On part du texte brut, on le
nettoie, on le transforme en nombres, on peut rééquilibrer les classes et réduire la
dimension, puis on entraîne un modèle réglé, et enfin on évalue. Deux choix
d'ingénierie : on nettoie une seule fois car cette étape n'apprend rien, ce qui rend
tout plus rapide, et l'équilibrage comme la réduction ne touchent que l'entraînement,
jamais le test, pour que nos résultats restent honnêtes.
\end{turn}

\begin{turn}{yann}{Slide 6 : Prétraitement}
Le prétraitement suit la séquence classique : minuscules, retrait de la ponctuation, des
URL et des chiffres, puis tokenisation et suppression des mots vides. Deux finesses à
retenir. On conserve les négations, car not good et good ont des sens opposés. Et la
lemmatisation regroupe les formes d'un mot, crashing et crash deviennent le même mot, ce
qui réduit le vocabulaire. Je passe la parole à Vénus pour la vectorisation.
\end{turn}

% ============================ VENUS ===================================
\section*{\textcolor{venus}{Partie 2 : Vénus BAKIKO (slides 7 à 12)}}

\begin{turn}{venus}{Slide 7 : Vectorisation}
Merci Yann. Les modèles ne comprennent que des nombres, il faut donc transformer le
texte. On compare quatre représentations demandées par le sujet. Bag-of-Words compte les
mots. TF-IDF fait pareil mais valorise les mots rares et discriminants. Word2Vec
représente chaque mot par un vecteur dense, et on en a deux versions, une entraînée sur
nos avis et une pré-entraînée sur Google News, qui est exactement le pre-trained Word2Vec
du sujet. Détail utile : Naive Bayes multinomial exige des valeurs positives, donc pour
les embeddings on passe à la version gaussienne.
\end{turn}

\begin{turn}{venus}{Slide 8 : Les cinq modèles}
Nous comparons cinq classifieurs, du plus simple au plus riche. Naive Bayes suppose les
mots indépendants, c'est faux mais très efficace, et ce sera notre gagnant. L'arbre de
décision est lisible mais surapprend, d'où l'élagage. La forêt aléatoire moyenne
beaucoup d'arbres. AdaBoost combine des souches, mais sur du texte creux il est faible.
La régression logistique est une référence solide avec régularisation. Tous sont réglés
par validation croisée en cinq blocs et recherche sur grille.
\end{turn}

\begin{turn}{venus}{Slide 9 : Les métriques}
Quatre mots à maîtriser. L'exactitude est la proportion globale de bonnes réponses, mais
elle trompe quand une classe domine. La précision répond à : quand le modèle annonce une
classe, a-t-il raison. Le rappel répond à : parmi les vrais cas, combien sont retrouvés.
Le F1 combine les deux. Et le macro-F1 fait la moyenne des F1 des trois classes à
égalité, donc rater la classe rare coûte cher. C'est pour cela qu'on règle sur le
macro-F1.
\end{turn}

\begin{turn}{venus}{Slide 10 : Résultats}
Voici la comparaison de toutes les combinaisons. Le meilleur est TF-IDF avec Naive Bayes,
à 0,559 de macro-F1. Le tableau donne le meilleur modèle pour chaque vectoriseur. Deux
observations : les cinq premiers sont tous des méthodes par comptage, ce qui est logique
sur des avis courts, et AdaBoost est dernier car ses souches ne voient qu'un mot à la
fois. Les scores sont serrés autour de 0,5, car tous butent sur la même difficulté.
\end{turn}

\begin{turn}{venus}{Slide 11 : Le piège de l'exactitude}
Ce slide est le cœur de notre message. Notre meilleur modèle affiche 84 pour cent
d'exactitude, ce qui paraît excellent. Mais la colonne neutre de la matrice de confusion
est entièrement vide : le modèle ne prédit jamais la classe neutre. Les 53 vrais avis
neutres sont tous classés en négatif ou positif. Le macro-F1 chute à 0,56 et révèle ce
problème. La leçon : sur des données déséquilibrées, l'exactitude ment, il faut regarder
le macro-F1 et le rappel par classe.
\end{turn}

\begin{turn}{venus}{Slide 12 : Le déséquilibre}
On attaque le déséquilibre. Sans traitement, le rappel de la classe neutre est exactement
zéro : la minorité est ignorée. SMOTE génère des exemples neutres synthétiques et fait
remonter ce rappel de zéro à 0,19. Le sous-échantillonnage va plus loin, à 0,47, mais en
jetant des données. Le rééquilibrage échange donc de l'exactitude contre une vraie
reconnaissance de la classe rare. Je laisse Faïçal pour les réglages et la conclusion.
\end{turn}

% ============================ FAICAL ==================================
\section*{\textcolor{faical}{Partie 3 : Faïçal DIELO (slides 13 à 18)}}

\begin{turn}{faical}{Slide 13 : Réduction de dimension}
Merci Vénus. La réduction de dimension. Notre matrice TF-IDF a environ 4\,600 colonnes.
On utilise TruncatedSVD, l'équivalent de la PCA adapté aux données creuses. Le résultat
est net : avec seulement 300 composantes, soit 6,5 pour cent des colonnes, on conserve
quasiment toute la performance. C'est un excellent compromis vitesse mémoire, et c'est
même indispensable avant un modèle dense comme le gradient boosting.
\end{turn}

\begin{turn}{faical}{Slide 14 : Élagage de l'arbre}
L'élagage de l'arbre de décision. Un arbre laissé libre construit 1\,857 nœuds et
mémorise les données : 0,98 en entraînement mais 0,76 en test, donc il surapprend. En
élaguant, on réduit l'arbre à 69 nœuds et la performance de test monte à 0,79, tout en
réduisant l'écart entre entraînement et test. Si on élague trop, vers 29 nœuds, l'arbre
sous-apprend. C'est l'illustration concrète du compromis biais variance.
\end{turn}

\begin{turn}{faical}{Slide 15 : Arrêt précoce}
L'arrêt précoce pour le gradient boosting, un modèle qui ajoute des arbres un par un. On
autorise jusqu'à 500 arbres, mais on surveille une part de validation et on arrête dès
que la performance ne progresse plus. Ici l'arrêt survient à 88 arbres sur 500, soit
environ six fois moins de calcul, sans perte de qualité, et une protection contre le
surapprentissage.
\end{turn}

\begin{turn}{faical}{Slide 16 : Conclusion}
Pour conclure. Nous avons construit une chaîne complète et reproductible sur de vrais
avis Facebook, comparé quatre représentations du texte et cinq modèles, tous réglés par
validation croisée. Le meilleur est TF-IDF avec Naive Bayes à 0,559 de macro-F1. La leçon
principale est méthodologique : sur un jeu déséquilibré, l'exactitude trompe, et c'est le
macro-F1 qui dit la vérité. SMOTE est ce qui permet à la classe minoritaire d'exister.
\end{turn}

\begin{turn}{faical}{Slide 17 : Limites et perspectives}
Soyons honnêtes sur les limites. Nos étiquettes viennent des étoiles, c'est un signal
approximatif. La classe neutre est minuscule et ambiguë. Et le texte est court, bruité et
parfois dans une autre langue. Côté perspectives, on peut activer BERT grâce au script
que nous fournissons, essayer la pondération des classes et des seuils ajustés, et
collecter de vraies issues GitHub ou JIRA.
\end{turn}

\begin{turn}{faical}{Slide 18 : Remerciements}
Merci de votre attention. Nous sommes maintenant prêts à répondre à vos questions. Pour
les questions sur les métriques, pensez au macro-F1 et au rappel par classe. Pour celles
sur le déséquilibre, pensez à SMOTE.
\end{turn}

\vspace{0.3cm}
{\color{accent}\rule{\linewidth}{1.2pt}}
\begin{center}\small\textcolor{primary}{Bonne présentation à tous les trois.}\end{center}

\end{document}
